\chapter{圈群、簇代数、二聚体与可积系统}
\label{chap:01}

\noindent Vladimir V. Fock、Andrey Marshakov\footnote{我们感谢奥胡斯大学 QGM、马克斯·普朗克研究所
与波恩 Haussdorf 数学研究所——本工作的重要部分是在这些地方完成的。我们感谢审稿人仔细阅读
手稿并改正了许多错误。V.F. 的工作部分受 ANR GTAA 与 ANR ETTT 基金资助；A.M. 的工作部分受
俄罗斯国立高等经济学院（NRU HSE）科研基金 13-05-0006、俄德联合项目 RFBR 12-02-92108、
科学学派支持计划（NSh-3349.2012.2）以及俄罗斯教育部合同 8207 资助。}

本文的主旨是证明：对于一大类可积模型，两种先验地不同的构造与描述方法是等价的，从而为
研究它们提供统一的进路。第一种方法 [24] 是熟知的：相空间取为 Kac--Moody 李群的双
Bruhat 胞腔之商，Poisson 结构由经典 $r$-矩阵给出，运动积分就是 $\Ad$-不变函数。第二种
方法由 A.~Goncharov 与 R.~Kenyon 新近在 [11] 中提出，源自统计物理中二维环面上二部图上的
二聚体模型的研究。我们将证明：后一类可积系统实际上是前者的特例，对应
$\hat A_{N-1}$ 型仿射群。

这类可积系统最著名的例子 [6] 是相对论 Toda 链，由 S.~Ruijsenaars [25] 发现，
Y.~Suris [28] 等许多人研究过；它在某个极限下给出通常的 Toda 链，用我们的语言说，这个
极限就是从李群过渡到李代数。另一个已知例子是五角星映射——射影平面中折线空间上的一个
离散可积系统，由 R.~Schwartz [26] 发现，他与 V.~Ovsienko、S.~Tabachnikov [22, 23]
研究过，近来 M.~Glick [10]，B.~Khesin 与 F.~Soloviev [17, 18, 27]，以及
M.~Shapiro、M.~Gekhtman、A.~Vainshtein 与 S.~Tabachnikov [9] 也研究过。另一方面，
V.~Ovsienko 与 S.~Tabachnikov 在 [22] 中证明：五角星映射的离散流在连续极限下给出
Boussinesq 流。这一观察在 [17] 中对其他维数作了推广，它提示 [22] 提出的技巧也可用于
研究连续可积系统，例如 $n$-KdV 梯队。

\section{引言}

\subsection{Poisson--Lie 群上的可积系统}

我们的出发点是：在 Poisson--Lie 群上，若 Poisson 括号由经典 $r$-矩阵定义，则
$\Ad$-不变函数彼此 Poisson 交换。对有限维单群，线性无关的 $\Ad$-不变函数的个数等于
群的秩，因此相应的可积系统只能出现在维数不超过秩的二倍的辛叶上，见 [6]。然而对仿射群，
线性无关的 $\Ad$-不变函数有无穷多个，而所有 Poisson 子流形仍是有限维的，于是这样得到的
可积系统集合要大得多。这类可积系统可以在任意仿射 Poisson--Lie 群 $\hat G$ 上构造，
但这里我们只需要 $\hat A$ 型、中心平凡的那些群上的系统。这个群可以实行为单变量 $\lambda$
的矩阵值 Laurent 多项式 $A(\lambda)$ 构成的群，其行列式是非零常数，并且按乘非零常数取商。

对给定的 $A(\lambda)$，集合
\[
\{(\lambda,\mu)\mid \det(A(\lambda)-\mu)=\sum_{ij}H_{ij}\lambda^i\mu^j=0\}
\]
是 $(\C^\times)^2$ 中的一条代数曲线，称为\emph{谱曲线}（spectral curve），它嵌入环面
$(\C^\times)^2$ 中。这条曲线若按环面自同构取商，则是共轭类的不变量。谱曲线连带一个
线丛，即由 $A(\lambda)-\mu\cdot\Id$ 的核给出的线丛。从群 $\hat G$ 到曲线空间的映射称为
\emph{作用映射}（action map）；若在曲线空间上取平凡 Poisson 括号，则它是 Poisson 映射。
$H_{ij}$ 本身并未良好定义，因为作变换 $H_{ij}\to H_{ij}\alpha^i\beta^j\gamma$ 后对应
同一条曲线；不过可以利用这一自由度，使任意三个非零系数等于 $1$。加上这个条件后，
$H_{ij}$ 就良定义，并且彼此 Poisson 交换。到（曲线，其上的线丛）这一对儿的映射称为
\emph{作用--角映射}（action-angle map）。由彼此 Poisson 交换的运动积分（或哈密顿量）
生成的流，归结为线丛沿谱曲线的 Jacobian 作常值流动。§1.2 将更详细地描述这类可积系统。

圈群 $\hat G$ 本身没有簇代数簇结构。不过，它作为 Poisson 子流形嵌入到中心余扩张
$\hat G^\sharp$ 中；后者有标准分解，即分解为 Poisson 子流形 $(\hat G^\sharp)_u$ 的
不交并，这些子流形在 [19] 中称为\emph{双 Bruhat 胞腔}（double Bruhat cell），它们
已经是簇代数簇。这些胞腔由余扩张 $(\tilde W\times\tilde W)^\sharp$ 的元素 $u$ 编号；该
余扩张是 $\hat G$ 的外尔群 $\tilde W$ 的平方被 Dynkin 图的自同构群所作的扩张，对
$\hat A_{N-1}$ 系列这个自同构群是循环群。

$\hat G^\sharp$ 的 Bruhat 胞腔与 $\hat G$ 的交，再商掉有限维 Cartan 子群 $H$ 的共轭
作用（我们也称之为双 Bruhat 胞腔，记作 $\hat G_u$），就是我们的可积系统的相空间
$\hat G_u/\Ad H$。这样的相空间维数（在仿射群情形）比 $u$ 的长度小 $1$。把 $u$ 写成
标准生成元的一个约化词，便可定义\emph{簇坐标} $x=\{x_f\}$：它们由词的字母（余扩张的
生成元除外）编号，并满足 $\prod x_f=1$。对一个给定的 Bruhat 胞腔，Laurent 多项式
$\det(A(x,\lambda,\mu))=\sum H_{ij}(x)\lambda^i\mu^j$ 有固定的 Newton 多边形 $\Delta$；
在该胞腔上不恒为零的、彼此 Poisson 交换的哈密顿量的个数，比 $\Delta$ 的整点个数少 $3$。
为使这个系统可积，这些运动积分必须相互独立，而且其个数应在给定胞腔维数与 Poisson 括号
秩的前提下达到最大可能值 $(\dim\hat G_u+\operatorname{corank}\hat G_u)/2$。在对应于
「$u$ 在其共轭类中长度最小」的双 Bruhat 胞腔上（这样的元素也称为循环不可约的），这个
条件成立。

沿用 [4]，我们在 §1.4 对群 $\hat G^\sharp$ 的双 Bruhat 胞腔及其按 Cartan 子群共轭的
商引入簇坐标（簇种子）；先在 §1.3 对单群描述，然后推广到仿射情形。给定与某个约化词
对应的一组坐标 $x$，我们把相应的矩阵多项式 $A(x,\lambda)$ 构造为初等矩阵的乘积，其中
每个因子或是常数，或只依赖单个坐标 $x_f$。然后，我们用这些簇坐标表述可积系统。同一个
词 $u$ 的不同分解所对应的坐标之间，由簇变换相联系。

\subsection{Goncharov--Kenyon 可积系统}

回忆：\emph{簇代数簇}（cluster variety）是一个代数簇，它被一族与代数环面 $(\C^\times)^N$
同构的坐标卡覆盖，转移函数是称为\emph{变异}（mutation）的特殊双有理变换的复合（见
附录 1.9）。Goncharov--Kenyon（GK）进路提出一个可积系统，它与二维环面上某个图上的
二聚体模型相联系。结果表明，这个可积系统结构与簇代数簇结构相容，从而带来对偶、离散群
作用、量子化、热带极限以及许多其他吸引人的性质。此外，彼此 Poisson 交换的运动积分可以
取为簇函数。许多这样的系统容许一个 Abel 群，由离散簇变换组成，并与可积流交换。

GK 可积系统构造的大致框架如下（详见 §1.5）。出发点是闭曲面 $\Sigma$ 上的一个二部图，
它满足某些极小性与非平凡性条件（沿用 [11]，我们设曲面 $\Sigma$ 的亏格为 $1$）。构造的
许多方面可以推广到更高亏格的曲面，但这一推广留待另文。

考虑图 $\Gamma$ 上取值于某个乘法群的离散联络空间（具体取为非零复数乘法群 $\C^\times$）。
二部图的每条边都有典范定向，例如从白顶点指向黑顶点，于是可以把该空间解释为 $1$-上链的
乘法群 $C^1(\Gamma)$，它无非是对应于 $\Gamma$ 每条边的 $\C^\times$ 因子的乘积。这个
空间按离散规范变换取商，可以解释为上同调群 $H^1(\Gamma)$。该群是上边缘群 $B^2(\Sigma)$
被环面的上同调群 $H^1(\Sigma)=\C^\times\times\C^\times$ 所作的扩张。$B^2(\Sigma)$ 的
元素是 $\Sigma$ 各面上的一族数，其乘积等于 $1$。微分
$\partial:H^1(\Gamma)\to B^2(\Sigma)$（由绕面的单值给出）在 $H^1(\Gamma)$ 上定义了
$B^2(\Sigma)$ 上的主 $H^1(\Gamma)$-丛结构。

对每个离散联络 $\mathbb{A}=\{A_e\}\in C^1(\Gamma)$，并在环面 $\Sigma$ 上取定一个自旋
结构，我们在 §1.5.5 定义一个\emph{离散 Dirac 算子}
$D(\mathbb{A}):C_B\to C_W$，它把黑顶点上的函数送到白顶点上的函数。该算子在
$C^1(\Gamma)$ 的一个子簇上退化，这个子簇是行列式 $\det D(\mathbb{A})$ 的零点集。
该子簇在规范变换下不变，因而定义了 $H^1(\Gamma)$ 中的一个子簇。这个零点集与
$B^2(\Sigma)$ 中点 $x=\{x_f\}$ 上的纤维之交，给出了一条代数曲线，其上的线丛由
$D(\mathbb{A})$ 的核给出。注意 $\det D(\mathbb{A})$ 是白、黑顶点的完美匹配上的单项式
之和，二聚体组态正是在此处进入图 $\Gamma$ 的构造。自旋结构使我们能控制这些单项式的符号。
于是这个构造定义了可积系统相空间 $B^2(\Sigma)$ 到平面代数曲线空间的作用映射，以及到
（平面曲线，其上的线丛）对儿的作用--角映射，后者是双有理同构。

在坐标中，作用映射如下。取丛 $\partial$ 的一个平凡化，这相当于取定同构
$H^1(\Gamma)=B^2(\Sigma)\times H^1(\Sigma)$，然后取 $H^1(\Gamma)$ 到 $C^1(\Gamma)$ 的
一个提升。在这些等同之下，我们可以把联络 $\mathbb{A}(x,\lambda)$ 联系到任意
$x\in B^1(\Sigma)$ 与 $\lambda=(\lambda,\mu)\in H^1(\Sigma)$。谱曲线由方程
$\det D(\mathbb{A}(x,\lambda))=\sum H_{ij}(x)\lambda^i\mu^j=0$ 定义；若把曲线按
$H^1(\Sigma)=\C^\times\times\C^\times$ 的自同构取商，则它不依赖所作的选择。系数
$H_{ij}$ 定义到变换 $H_{ij}\to H_{ij}\alpha^i\beta^j\gamma$ 为止，我们可以利用这个
自由度使其中三个等于 $1$。其余的系数给出 $B^2(\Sigma)$ 上自然 Poisson 结构下的一族
彼此 Poisson 交换的函数；若图满足某个极小性条件，这个 Poisson 结构是最大的。

空间 $B^2(\Sigma)$ 具有 [11] 中如下引入的典范对数常值 Poisson 括号。把图 $\Gamma$
嵌入曲面 $\Sigma$ 会在每个顶点处诱导边端的一个循环次序（厚图结构）。考虑与同一个二部图
$\Gamma$ 相对应的另一个曲面 $\tilde\Sigma$，但在白顶点处把循环次序改为相反，在黑顶点处
保持不变。由于 $\Gamma$ 嵌入 $\tilde\Sigma$，我们有映射 $H^1(\tilde\Sigma)\to B^2(\Sigma)$，
它是标准嵌入与上边缘算子之复合。空间 $H^1(\tilde\Sigma)$ 具有由 $\tilde\Sigma$ 上相交
指数给出的典范 Poisson 结构，而到 $B^2(\Sigma)$ 的映射在后一空间上诱导出 Poisson 括号。

这个括号可以通过乘法不变性延拓到空间 $B^2(\Sigma)$——即附在图的各面上的一族非零数
构成的空间——并定义一个簇种子，其反对称交换矩阵由该 Poisson 括号给出。在 [11] 中注意
到：图容许初等变换，称为\emph{蛛形变换}（spider move），只要相空间之间由一个簇变异相
联系，与这些图相应的可积系统就同构。在这种变换下可积系统的等价类由
$\det D(x,\lambda)$ 的 Newton 多边形编号，而独立运动积分的个数恰好是这些 Newton 多边形
严格内部的整点个数。

\subsection{两种进路之间的关系}

我们在 §1.6 断言：GK 可积系统与 Poisson--Lie 圈群 $\widehat{PGL}(N)$ 上的可积系统
相重合。这个同构不仅等同它们的相空间与交换流，还等同离散群作用与典范簇坐标。

在两个构造中，可积系统的谱曲线都由某个矩阵算子的退化条件给出（群论进路中是
$A(x,\lambda)-\mu\cdot\Id$，另一侧是 Dirac 算子 $D(\mathbb{A}(x,\lambda,\mu))$）。
这两个矩阵本身并不相同，但它们的行列式相同；粗略地说，两者的对应关系如下。任意矩阵 $D(\mathbb{A})$
的行列式都可以写成 Grassmann 积分：\hfil
\[
S_\Gamma(\mathbb{A})=\det D(\mathbb{A})
=\int \exp\Bigl(\sum \mathfrak{D}(\mathbb{A})^b_w\,\xi_b\eta^w\Bigr)
\prod_b d\xi_b\prod_w d\eta^w .
\]
因此，$\det D(\mathbb{A})$ 可以解释为某个格点费米子在背景规范场 $\mathbb{A}$ 中的
配分函数 $S_\Gamma(\mathbb{A})$。把环面剪成柱面，相当于把这个配分函数改写成演化算子
从一侧边界圆到另一侧边界圆的迹。这个演化算子由矩阵 $A(\lambda)$ 给出，后者作用在
$N$ 维空间的外代数上。进一步把这个柱面剪成一串更小的柱面，就可以把演化算子表示成
一系列初等步骤的乘积，每一步至多依赖一个变量 $x_f$，并且恰好与用来参数化双 Bruhat
胞腔的初等矩阵重合，从而建立了谱曲线之重合。

这个纲领的第一部分，是建立 $(\tilde W\times\tilde W)^\sharp$ 的生成元中的词（它们在第一
种进路中编号簇坐标系）与二部图（它们在第二种进路中编号坐标）之间的对应。此外，这个
对应应当把词的字母与相应二部图（画在环面 $\Sigma$ 上）的面等同起来，因为在两种情形中
这两组对象都对应簇坐标。

为此我们使用第三个组合对象，它由 Dylan Thurston 在未刊论文 [29] 中提出（在我们的
语境中 [11] 已经使用过），我们称之为 \emph{Thurston 图}，并在附录 1.12 中详述。
Thurston 图是曲面上的一族曲线的同痕类，这些曲线或是闭的，或是连接两个边界点，只有
三重点，并且补集的连通分支（面）染成白色与灰色，任意两个共享一段曲线的面颜色不同
（棋盘染色）。这样的图容许初等修改，称为 \emph{Thurston 变换}。正如 D.~Thurston 与
A.~Henriques 已经注意到的，每个 Thurston 图定义一个簇种子（簇流形上的一个坐标卡），
其簇变量附在白面上。Thurston 变换对应变异（从一个坐标卡过渡到另一个）。有了开曲面上的
Thurston 图，就可以尊重染色地把边界分支粘起来，从而得到带 Thurston 图的新曲面。

为了从元素 $u\in(\tilde W\times\tilde W)^\sharp$ 的约化分解构造 Thurston 图，我们首先把
$(\tilde W\times\tilde W)^\sharp$ 的每个生成元（余中心生成元除外）联系到一个柱面上的
Thurston 图，它只有一个三重点，且在每个边界圆上有 $N$ 段灰色（与白色）线段。然后，
按约化分解中生成元的次序把这些柱面粘起来，最后把所得柱面的两端带一个扭转粘起来，扭转
由余中心生成元 $\Lambda$ 的幂给出。

为了从 Thurston 图构造二部图，我们在每个三重点放一个黑顶点，在每个灰色面放一个白顶点。
然后从每个黑顶点向该顶点所接的三个灰色扇区内各画一条边，连到相应的白顶点。容易看出，
约化词的字母集合与 Thurston 图的白面集合有典范双射，而后者又与二部图的面双射。下一个
观察同样简单：这个双射诱导簇种子之间的双射，即坐标之间的 Poisson 括号相重合。

最后，我们需要证明方程 $\det(A(x,\lambda)-\mu)=0$ 与 $D(\mathbb{A}(x,\lambda,\mu))=0$
定义同一条曲线。为此，我们在 §1.5 把二部图上的格点费米子配分函数延拓到带边曲面。
这类曲面上的图容许第三类顶点：它们终止在边界上，可以连到白顶点，也可以连到黑顶点，
但彼此不相连。记这类顶点的集合为 $T$。Dirac 算子此时作用为
$D(\mathbb{A}):C_{B\cup T}\to C_{W\cup T}$；再引入额外的 Grassmann 变量
$\zeta=\{\zeta_t\mid t\in T\}$，定义
\[
S(\mathbb{A},\zeta)=\int \exp\Bigl(\sum \mathfrak{D}(\mathbb{A})^b_w\,\xi_b\eta^w
+\mathfrak{D}(\mathbb{A})^b_t\,\xi_b\zeta^t
+\mathfrak{D}(\mathbb{A})^t_w\,\zeta_t\eta^w\Bigr)
\prod_b d\xi_b\prod_w d\eta^w .
\]
把 $\Sigma$ 的两个边界分支沿二部图 $\Gamma$ 粘起来——粘法使终止顶点与终止顶点相接——
就得到粘合曲面 $\bar\Sigma$ 上的二部图 $\bar\Gamma$。$\Gamma$ 上的联络 $\mathbb{A}$ 诱导
$\bar\Gamma$ 上的联络 $\bar{\mathbb{A}}$：只需把粘合边两半的数相乘。若
$S_\Gamma(\mathbb{A},\zeta)$ 是 $\Gamma$ 的配分函数，则 $\bar\Gamma$ 的配分函数为
\[
S_{\bar\Gamma}(\zeta,\bar{\mathbb{A}})
=\int S_\Gamma(\zeta,\mathbb{A})\,
\exp\Bigl(\sum_t \zeta_t\zeta_{\sigma(t)}\Bigr)
\prod d\zeta_t\,d\zeta_{\sigma(t)},
\]
其中指标 $t$ 跑遍一侧的终止边，$\sigma$ 是把终止顶点送到与之粘合的顶点的映射，而
$\bar\lambda$ 是从 $\lambda$ 中去掉对应于被粘顶点的项后所得。

另一方面，注意任意 $N\times N$ 矩阵 $M$ 都可以联系到一个 $2N$ 个 Grassmann 变量
$\xi=\{\xi_i\}$ 与 $\eta=\{\eta_i\}$ 的函数 $S_M(\xi,\eta)=\exp(M^i_j\xi_i\eta^j)$。
矩阵乘积于是对应相应函数的卷积：
\[
S_{M_1M_2}(\xi,\eta)=\int S_{M_1}(\xi,\eta')S_{M_2}(\xi',\eta)\,
\exp\Bigl(\sum \xi'_i\eta'^i\Bigr)\prod d\xi'_i\,d\eta'^i .
\]
因此，若柱面上某个图的配分函数与某个矩阵的配分函数重合，则若干柱面粘合后的配分函数
对应矩阵的乘积。于是，为证明曲线重合，我们只需把环面剪成小柱面，并对每个小柱面验证
配分函数重合。我们在 §1.6 完成对应关系的证明，并在那里陈述主要结果。

§1.7 讨论可积系统中的离散流。§1.8 讨论可积系统的几个例子。特别地，§1.8.5 将证明：
R.~Schwartz [26] 发现的射影平面中多边形空间上的离散可积系统，可以实为本方案的一个
特例。这一类中更特殊的系统例子已在 [21] 中考虑过。

\section{可积系统与 \texorpdfstring{$r$}{r}-矩阵}

回忆与单李群或仿射李群上的经典 $r$-矩阵相关的可积系统的标准构造（例如见 [24]）。
我们设这些群是复的、中心为零。这类可积系统的相空间是 Poisson--Lie 群的某些 Poisson
子流形，而彼此交换的哈密顿量或运动积分由共轭不变函数给出。事实上，设 $G$ 是李群，
$\mathfrak{g}=\operatorname{Lie}(G)$ 是相应的李代数，$r\in\mathfrak{g}\otimes\mathfrak{g}$
是 Yang--Baxter 方程
\[
[r_{12},r_{13}]+[r_{13},r_{23}]+[r_{12},r_{23}]=0
\]
的解。这样的 $r$-矩阵在群 $G$ 上定义 Poisson 括号
\begin{equation}
\{g\otimes,g\}=-\frac12\,[r,\,g\otimes g],\tag{1.1}\label{eq:1.1}
\end{equation}
并且这个括号与群结构相容，意思是群乘法 $G\times G\to G$ 与取逆 $G\to G$ 都是 Poisson
映射。

从 (1.1) 可以直接看出：$G$ 上任意两个 $\Ad$-不变函数彼此 Poisson 交换。事实上，把
$r$-矩阵写成 $r=\sum_I v_I^{(1)}\otimes v_I^{(2)}$，其中所有 $v_I\in\mathfrak{g}$，
并用 $L_v$（分别地 $R_v$）记对应于 $v$ 的左（分别地右）不变向量场，则任意两个函数的
Poisson 括号为
\[
\{H_1,H_2\}=\sum_I\Bigl(L_{v_I^{(1)}}H_1\,L_{v_I^{(2)}}H_2
-R_{v_I^{(1)}}H_1\,R_{v_I^{(2)}}H_2\Bigr).
\]
由于任意 $\Ad$-不变函数 $H$ 满足 $L_vH=-R_vH$，两个这样的函数的括号为零。还要注意，
这个论证是局部的：即使函数不是定义在整个 $G$ 上，而是定义在 $G$ 的任意 Poisson
$\Ad$-不变子簇上，括号仍然为零。

我们把自己限制在单李群或仿射李群 $\hat G$ 的情形，此时 Yang--Baxter 方程有典范的
Drinfeld--Jimbo 解
\begin{equation}
r=\sum_{\alpha\in\Delta_+}d^\alpha e_\alpha\otimes e_{\bar\alpha}
+\frac12\sum_{i\in\Pi}d^i h_i\otimes h_i,\tag{1.2}\label{eq:1.2}
\end{equation}
其中 $\Delta_+$ 是正根集合，$d^\alpha=(\alpha,\alpha)/2$，$\Pi$ 是正单根集合，
$\bar\alpha$ 只是 $-\alpha$ 的另一种记法，而 $e_\alpha$ 与 $h_i$ 构成 $\mathfrak{g}$
的标准 Cartan--Weyl 基。为简化叙述，以下设群是单连通的，即 $(\alpha,\alpha)=2$，且
对所有根 $\alpha$ 有 $d^\alpha=1$。

显然，单群上存在 $\operatorname{rank}G$ 个独立的 $\Ad$-不变函数：一种选法是集合
$\{\mathcal{H}_i\}$，其中 $i\in\Pi$，$\mathcal{H}_i(g)=\operatorname{Tr}\pi_{\mu^i}(g)$，
而 $\pi_{\mu^i}$ 是 $G$ 的第 $i$ 个基本表示，其最高权 $(\mu^i,\alpha_j)=\delta^i_j$
与 $\alpha_i$ 对偶，$i\in\Pi$。这些函数在 $G$ 的秩至多为 $2\operatorname{rank}G$ 的
Poisson 子流形上定义可积系统（例如见 [6]）。对圈群 $\hat G$——以下我们把它理解为取值
于单群 $G$ 的 Laurent 多项式之群——独立 $\Ad$-不变函数的个数是无穷的，因为
$\operatorname{Tr}\pi_{\mu^i}(g)$ 的每个系数现在都是 $\Ad$-不变函数。于是圈群给出
大得多的可积模型集合。

在空间 $G/\Ad H$ 上（$H$ 是 $G$ 的 Cartan 子群），沿用 [12] 可以定义一个离散双有理
Poisson 变换 $\tau:G/\Ad H\to G/\Ad H$ 的作用，它保持双 Bruhat 胞腔与函数
$\mathcal{H}_i$。具体说，设 $g=g_+g_-$ 是 $g\in G/\Ad H$ 的 Gauss 分解，定义
$\tau(g)$ 为乘积 $g_-g_+$。（Gauss 分解在 $g$ 上并非唯一确定，因为 Cartan 部分既可
附到上三角部分，也可附到下三角部分，或分拆到两者；但在商 $G/\Ad H$ 上，$\tau$ 的作用
仍然是良定义的。）在 §1.7 我们将说明如何把这个变换推广到圈群。
